\begin{table}[htbp]\centering
\begin{threeparttable}
\caption{SHs and a VT that are reached by divergent sequences but by at most one main-type sequence (paralog names).}
\label{tab:paralog}
\footnotesize\setlength{\tabcolsep}{4pt}
\begin{tabular}{L{0.165\linewidth}L{0.20\linewidth}L{0.325\linewidth}L{0.06\linewidth}L{0.13\linewidth}}
\toprule
Taxon code & Taxon of the best hit & Divergent sequences & Main$^{a}$ & Identity$^{b}$\\
\midrule
SH1337353.10FU & \emph{Rhizoglomus irregulare} & \emph{Rhizoglomus irregulare} 10 (7 cultures) & 1 & 100.0\%\\
SH1364117.10FU & Glomeromycota sp. & \emph{Rhizoglomus vesiculiferum} 7 (2 cultures) & 0 & 98.5--99.8\%\\
SH1081549.10FU & \emph{Rhizoglomus} sp. & \emph{Rhizoglomus proliferum} 4 (1 culture) & 0 & 99.3--99.8\%\\
SH1081557.10FU & \emph{Rhizoglomus intraradices} & \emph{Rhizoglomus intraradices} 4 (3 cultures) & 0 & 100.0\%\\
SH1076762.10FU & \emph{Gigaspora margarita} & \emph{Gigaspora rosea} 2 (1 culture) & 1 & 99.8\%\\
SH1337340.10FU & \emph{Rhizoglomus} sp. & \emph{Rhizoglomus} sp. 1; \emph{Rhizoglomus irregulare} 1 (2 cultures) & 0 & 100.0\%\\
SH1337458.10FU & \emph{Rhizoglomus irregulare} & \emph{Rhizoglomus silesianum} 2 (1 culture) & 0 & 100.0\%\\
SH1389224.10FU & Glomeraceae sp. & \emph{Septoglomus viscosum} 2 (1 culture) & 0 & 100.0\%\\
VTX00409 & environmental (MaarjAM) & \emph{Septoglomus viscosum} 3 (2 cultures) & 0 & 98.5--99.4\%\\
\bottomrule
\end{tabular}
\begin{tablenotes}\footnotesize
\item $^{a}$\,Main-type sequences of any culture with that name. $^{b}$\,BLAST identity of the divergent sequences to the best hit. Taxa are those of the best UNITE record (MaarjAM's for the VT). 6 further names take one divergent
sequence each (SH1076712, SH1196937, SH1337247, SH1337463, SH1337480, SH1364178). 
\end{tablenotes}
\end{threeparttable}
\end{table}
